% sections/02_setup.tex -- Setup: model, data, teacher, training and diagnostics.
% Sources of truth: notes/model.md (model, init, dynamics, conservation laws) and
% notes/diagnostics.md (events, plateaus, regime rules). Keep this section in sync with them.
\section{Setup: model, data, teacher, training and diagnostics}\label{sec:setup}

This section fixes the model studied throughout and the measurements used to classify training
regimes. Every experiment uses these definitions; deviations are stated where they occur.

% ===========================================================================================
\subsection{Tree tensor network regression}\label{sec:setup-model}

\paragraph{Inputs.}
An input $x$ consists of $\Ns$ sites. Site $i\in\{1,\dots,\Ns\}$ carries a local feature vector
$\feat{i}(x)\in\R^{\pd[i]}$ (by default $\pd[i]=d$). Unless stated otherwise
$\feat{i}(x)=\xs{i}$ with $\xs{i}\sim\Normal(0,\Sigi{i})$ independent across sites, and the
inputs are \emph{whitened}, $\Sigi{i}=\Id$. We write
\begin{equation*}
  \Feat(x)=\feat{1}(x)\otimes\dots\otimes\feat{\Ns}(x)\in\R^{\pd[1]\times\dots\times\pd[\Ns]},
  \qquad
  \Sig=\Sigi{1}\otimes\dots\otimes\Sigi{\Ns}.
\end{equation*}

\paragraph{Trees and cores.}
A tree tensor network (TTN) is specified by a rooted tree $\tree$ with node set $\nodes$,
$\abs{\nodes}=\Kc$, and an assignment of every site to exactly one node. For a node $v$,
$\sites{v}$ is the (possibly empty) set of sites attached to $v$, $\ch{v}$ the set of its children
and $\pa{v}$ its parent. Every non-root node $v$ is joined to its parent by the edge (bond)
$\edge{v}$ of \emph{bond dimension} $\bd{v}$; for the root we set $\bd{\vroot}=1$. Node $v$ carries a
\emph{core} $\G{v}$, a tensor with one \emph{physical leg} of dimension $\pd[i]$ for every
$i\in\sites{v}$, one leg of dimension $\bd{c}$ for every child $c\in\ch{v}$, and one leg of
dimension $\bd{v}$ for its parent bond (for the root, the scalar output). The legs other than the
parent leg are the \emph{child legs}; their dimensions multiply to the \emph{fan-in}
\begin{equation*}
  \fanin{v}=\prod_{i\in\sites{v}}\pd[i]\prod_{c\in\ch{v}}\bd{c},
\end{equation*}
so $\G{v}$ can be viewed as an $\fanin{v}\times\bd{v}$ matrix. The parameters are
$\params=(\G{v})_{v\in\nodes}$, with $\sum_v\fanin{v}\bd{v}$ entries in total.

\paragraph{Model.}
The tensor $\W(\params)\in\R^{\pd[1]\times\dots\times\pd[\Ns]}$ is the contraction of all cores
over all bonds, and the model is
\begin{equation}\label{eq:model}
  f_\params(x)=\ip{\W(\params)}{\Feat(x)}.
\end{equation}
Equivalently, $f_\params$ is computed by passing messages $m_v(x)\in\R^{\bd{v}}$ from the leaves to
the root,
\begin{equation}\label{eq:messages}
  m_v(x)_a=\sum_{(j_i),(b_c)}\G{v}\big[(j_i)_{i\in\sites{v}},(b_c)_{c\in\ch{v}},a\big]
  \prod_{i\in\sites{v}}\feat{i}(x)_{j_i}\prod_{c\in\ch{v}}m_c(x)_{b_c},
  \qquad
  f_\params(x)=m_{\vroot}(x).
\end{equation}
Two structural facts drive the dynamics. The model is multilinear in the site features, and $\W$ is
multilinear in the cores (each core appears once), so $\params\mapsto\W(\params)$ is a
homogeneous polynomial of degree $\Kc$: rescaling the cores, $\G{v}\mapsto c_v\G{v}$, rescales $\W$
by $\prod_v c_v$.

\paragraph{Bipartitions and bond spectra.}
Let $\below{v}$ be the set of sites in the subtree rooted at $v$. The edge $\edge{v}$ splits the
sites into $\below{v}$ and its complement, and $\mat{\W}{\edge{v}}$ denotes the matricization of
$\W$ with the sites in $\below{v}$ as row indices. Contracting the subtree below $v$ and the rest of
the network separately gives a factorization $\mat{\W}{\edge{v}}=\Sub{v}C_v\T$ with
$\Sub{v}\in\R^{(\prod_{i\in\below{v}}\pd[i])\times\bd{v}}$, so
$\rank\mat{\W}{\edge{v}}\le\bd{v}$ \citep{hackbusch2009new,grasedyck2010hierarchical}. The
\emph{bond spectrum} (Schmidt values) of an edge $e=\edge{v}$ is the list
$\bs{e}{1}\ge\bs{e}{2}\ge\dots\ge\bs{e}{\bd{v}}\ge0$ of the $\bd{v}$ largest singular values of
$\mat{\W}{e}$. All nonzero singular values are among them, so $\sum_k(\bs{e}{k})^2=\normF{\W}^2$ on
every edge. With anisotropic inputs we use function-space spectra, the singular values of
$\mat{(\Sig^{1/2}\W)}{e}$ with $\Sig^{1/2}=\Sigi{1}^{1/2}\otimes\dots\otimes\Sigi{\Ns}^{1/2}$.
Bond spectra are computed from the cores by QR sweeps, without forming $\W$, to about $10^{-15}$
relative accuracy.

\paragraph{Why this model.}
It is the simplest genuine TTN. (i)~With at least three sites and an internal bond, $\W$ is a
genuine tensor rather than a matrix; a deep linear network is the special case of a chain with
physical legs only at its two ends. (ii)~Whitened inputs reduce the population loss to a Frobenius
distance (\cref{sec:setup-loss}), the exact analogue of the setting in which \citet{saxe2014exact}
solved deep linear networks, so that theory provides baselines. (iii)~It is the regression version
of the tensor-network models used in machine learning: the matrix product state (MPS) models of
\citet{stoudenmire2016supervised} have the same form $\ip{\W}{\bigotimes_i\feat{i}(x)}$, with
non-Gaussian local features. (iv)~With whitened inputs it is gradient flow on a fully observed
hierarchical tensor factorization, whose implicit regularization was studied by
\citet{razin2022implicit}.
% TODO(lit): once notes/literature.md is final, check that the factorization of Razin et al.
% (2022) matches our HT tree exactly (leaf cores, local components) and adjust (iv).

% ===========================================================================================
\subsection{Tree topologies}\label{sec:setup-trees}

\Cref{tab:trees} lists the built-in topologies and \cref{fig:trees} draws them. A binary tree and
an MPS with the same $\Ns$ are the two extremes of tree shape at nearly equal $\Kc$ ($\Ns-1$ versus
$\Ns$) and are the default pair for saddle-to-saddle studies. The hierarchical Tucker (HT) format
adds one leaf core per site; when the leaf bond dimension equals $d$ this adds no expressivity, but
it changes $\Kc$ and the number of cores that must change to learn a component. Any parent array
with a site assignment defines a custom tree. Unless stated otherwise every edge has the same nominal
bond dimension $\chi$, capped at the largest rank its bipartition allows,
\begin{equation}\label{eq:bond-cap}
  \bd{v}=\min\Big(\chi,\ \prod_{i\in\below{v}}\pd[i],\ \prod_{i\notin\below{v}}\pd[i]\Big).
\end{equation}

\begin{table}[t]
  \centering
  \caption{Built-in tree topologies for $\Ns$ sites ($\Kc$ cores). Depth is the number of edges on
    the longest root-to-leaf path. Binary-tree values are for $\Ns=2^m$; for other $\Ns$ the sites are
    split recursively into halves of sizes $\lceil\Ns/2\rceil$ and $\lfloor\Ns/2\rfloor$ until
    groups of at most two sites remain, each of which becomes a bottom core.}
  \label{tab:trees}
  \begin{tabularx}{\textwidth}{@{}lXcc@{}}
    \toprule
    name & structure & $\Kc$ & depth \\
    \midrule
    binary & balanced binary tree; each bottom core carries two adjacent sites, internal cores
             carry none & $\Ns-1$ & $\log_2\Ns-1$ \\
    MPS    & chain (caterpillar), the most unbalanced tree: core $k$ carries site $k$ and its parent
             is core $k+1$, so the root is core $\Ns$ & $\Ns$ & $\Ns-1$ \\
    HT     & hierarchical Tucker: balanced binary tree over $\Ns$ leaf cores, one site per leaf
             core; internal cores carry none & $2\Ns-1$ & $\lceil\log_2\Ns\rceil$ \\
    custom & any parent array and site assignment & -- & -- \\
    \bottomrule
  \end{tabularx}
\end{table}

% tikz/trees.tex -- figure: the built-in tree topologies (input from sections/02_setup.tex).
% Node layout follows ttn/tree.py: binary_tree, mps_tree (root = last core), ht_tree.
% Sites are numbered 1..N here (0..N-1 in the code).
\tikzset{
  ttn/core/.style={circle, draw=MidnightBlue!75!black, fill=MidnightBlue!12, line width=0.6pt,
                   minimum size=3.6mm, inner sep=0pt},
  ttn/root/.style={ttn/core, fill=MidnightBlue!60},
  ttn/bond/.style={line width=1.2pt, black!80},
  ttn/phys/.style={line width=0.6pt, black!55},
  ttn/site/.style={font=\scriptsize, inner sep=1pt, text=black!75},
}
\begin{figure}[t]
  \centering
  \begin{subfigure}[b]{0.37\textwidth}
    \centering
    \begin{tikzpicture}[x=6mm, y=8.5mm]
      \foreach \j in {1,...,4} \node[ttn/core] (b\j) at ({2*\j-1}, 0) {};
      \node[ttn/core] (m1) at (2, 1) {};
      \node[ttn/core] (m2) at (6, 1) {};
      \node[ttn/root] (r) at (4, 2) {};
      \draw[ttn/bond] (b1) -- (m1) -- (b2)  (b3) -- (m2) -- (b4)  (m1) -- (r) -- (m2);
      \foreach \j in {1,...,4} {
        \pgfmathtruncatemacro{\sa}{2*\j-1}
        \pgfmathtruncatemacro{\sb}{2*\j}
        \draw[ttn/phys] (b\j) -- ++(-0.45,-0.85) node[ttn/site, below] {\sa};
        \draw[ttn/phys] (b\j) -- ++(0.45,-0.85) node[ttn/site, below] {\sb};
      }
    \end{tikzpicture}
    \caption{binary, $\Ns=8$, $\Kc=7$}
    \label{fig:trees-binary}
  \end{subfigure}%
  \hfill
  \begin{subfigure}[b]{0.37\textwidth}
    \centering
    \begin{tikzpicture}[x=6mm, y=8.5mm]
      \foreach \k in {1,...,7} \node[ttn/core] (c\k) at (\k, 0) {};
      \node[ttn/root] (c8) at (8, 0) {};
      \foreach \k [evaluate=\k as \kn using int(\k+1)] in {1,...,7}
        \draw[ttn/bond] (c\k) -- (c\kn);
      \foreach \k in {1,...,8}
        \draw[ttn/phys] (c\k) -- ++(0,-0.85) node[ttn/site, below] {\k};
      \node[font=\scriptsize, above=1.5pt] at ($(c4)!0.5!(c5)$) {$\bd{4}$};
    \end{tikzpicture}
    \caption{MPS, $\Ns=8$, $\Kc=8$}
    \label{fig:trees-mps}
  \end{subfigure}%
  \hfill
  \begin{subfigure}[b]{0.22\textwidth}
    \centering
    \begin{tikzpicture}[x=7mm, y=8.5mm]
      \foreach \i in {1,...,4} \node[ttn/core] (l\i) at (\i, 0) {};
      \node[ttn/core] (h1) at (1.5, 1) {};
      \node[ttn/core] (h2) at (3.5, 1) {};
      \node[ttn/root] (r) at (2.5, 2) {};
      \draw[ttn/bond] (l1) -- (h1) -- (l2)  (l3) -- (h2) -- (l4)  (h1) -- (r) -- (h2);
      \foreach \i in {1,...,4}
        \draw[ttn/phys] (l\i) -- ++(0,-0.85) node[ttn/site, below] {\i};
    \end{tikzpicture}
    \caption{HT, $\Ns=4$, $\Kc=7$}
    \label{fig:trees-ht}
  \end{subfigure}
  \caption{The built-in tree topologies (\cref{tab:trees}). Circles are cores $\G{v}$; the shaded
    core is the root, whose parent leg (dimension one, the scalar output) is not drawn. Thick lines
    are bonds: the edge $\edge{v}$ between core $v$ and its parent has dimension $\bd{v}$ (one label
    shown in (b), where cores are numbered $1,\dots,8$ from the left). Thin dangling lines are
    physical legs, labelled by their site $i$ (dimension $\pd[i]$).}
  \label{fig:trees}
\end{figure}


% ===========================================================================================
\subsection{Teacher, data and loss}\label{sec:setup-loss}

\paragraph{Teachers.}
The target is realizable, $f^*(x)=\ip{\Wt}{\Feat(x)}$. A \emph{CP teacher} has
\begin{equation}\label{eq:cp-teacher}
  \Wt=\sum_{r=1}^{\ncomp}\tw{r}\U{r},\qquad
  \U{r}=\fac{r}{1}\otimes\dots\otimes\fac{r}{\Ns},\qquad
  \norm{\fac{r}{i}}=1,\qquad \tw{1}\ge\dots\ge\tw{\ncomp}>0.
\end{equation}
Its inner products with a TTN do not depend on the TTN's tree (see \eqref{eq:cp-inner}), and the
following special cases have exactly known bond spectra.
\begin{itemize}
  \item \emph{Odeco} (orthogonally decomposable): $\ip{\fac{r}{i}}{\fac{r'}{i}}=\delta_{rr'}$ at every
    site, which requires $\ncomp\le\min_i\pd[i]$. By default the factors at each site are columns of
    an independent Haar-random orthogonal matrix.
  \item \emph{Pattern teachers}: each site has an orthonormal basis $(\beta^{(i)}_1,\dots,
    \beta^{(i)}_{\pd[i]})$ and component $r$ uses $\fac{r}{i}=\beta^{(i)}_{\pi_r(i)}$ for a pattern
    $\pi_r\colon\{1,\dots,\Ns\}\to\{1,\dots,d\}$. Odeco teachers are the case where the patterns
    differ at every site; with $\pi_1=(1,1,1,1)$ and $\pi_2=(2,2,1,1)$ the second component differs from
    the first only on the sites $D=\{1,2\}$ (a \emph{local} component).
  \item \emph{Random CP}: $\fac{r}{i}=g/\norm{g}$ with independent $g\sim\Normal(0,\Id)$, so the
    components are generically non-orthogonal.
\end{itemize}
A \emph{TTN teacher} gives $\Wt$ as a TTN on the student's tree or on another tree; a random TTN
teacher is a fan-in initialization (\cref{sec:setup-init}) with $\isc=1$, rescaled so that
$\normF{\Wt}=1$. The overlap of the student with component $r$ of a CP teacher is
\begin{equation}\label{eq:overlap}
  \ov{r}=\ipS{\W}{\U{r}}/\tw{r};
\end{equation}
for an odeco teacher with whitened inputs, $\ov{r}=0$ at $\W=0$ and $\ov{r}=1$ at $\W=\Wt$.

\begin{proposition}[Teacher bond spectra]\label{prop:teacher-spectra}
  \proven{} Let $e$ be an edge whose two sides both carry sites (true for every edge of the trees of
  \cref{tab:trees}). (a)~For an odeco teacher the bond spectrum of $\Wt$ on $e$ is
  $(\tw{1},\dots,\tw{\ncomp})$, for every tree. (b)~For a two-component pattern teacher whose
  components differ exactly on the site set $D$, the bond rank of $\Wt$ on $e$ is $2$ if $e$ lies in
  the smallest subtree of $\tree$ containing the nodes that carry the sites of $D$, and $1$ otherwise.
\end{proposition}
\begin{proof}
  Let $A$ be the set of sites on one side of $e$ and write $x_r=\bigotimes_{i\in A}\fac{r}{i}$,
  $y_r=\bigotimes_{i\notin A}\fac{r}{i}$, so that $\mat{(\Wt)}{e}=\sum_r\tw{r}x_ry_r\T$. (a)~Both
  $(x_r)_r$ and $(y_r)_r$ are orthonormal families, so this is a singular value decomposition.
  (b)~If $D\subseteq A$ or $D\cap A=\emptyset$, then $y_1=y_2$ or $x_1=x_2$ and the rank is one.
  Otherwise the two components differ (by orthogonal basis vectors) on a site of $D\cap A$ and on a
  site of $D\setminus A$, so $x_1\perp x_2$, $y_1\perp y_2$ and the rank is two. Finally, $e$
  separates two sites of $D$ exactly when it lies on the tree path between the nodes carrying them,
  that is, when it belongs to the smallest subtree containing those nodes.
\end{proof}
The number of cores in that subtree is the quantity $k$ of the locality hypothesis tested in
\cref{sec:s2s}.

\paragraph{Population loss.}
Training minimizes the population square loss
\begin{equation}\label{eq:loss}
  \Loss(\params)=\tfrac12\,\E_x\big[(f_\params(x)-f^*(x))^2\big]
  =\tfrac12\normS{\W(\params)-\Wt}^2,
  \qquad \normS{A}^2=\ip{A}{\Sig A}.
\end{equation}
The second equality holds because, by independence across sites,
$\E[\Feat\otimes\Feat]=\bigotimes_i\E[\feat{i}\feat{i}\T]=\Sig$ (as an operator on
$\R^{\pd[1]\times\dots\times\pd[\Ns]}$). With whitened inputs $\Loss$ is half the squared Frobenius
distance. Only second moments enter: every local feature map with $\E[\feat{i}\feat{i}\T]=\Id$ gives
the same loss and the same dynamics (for example Hermite features of a scalar input, or
$(1,\xi_i)$ with $\xi_i\sim\Normal(0,1)$), and a non-orthonormal feature map is equivalent to
$\Sigi{i}=\E[\feat{i}\feat{i}\T]\ne\Id$.

\paragraph{Closed form.}
Expanding \eqref{eq:loss},
\begin{equation}\label{eq:loss-expanded}
  \Loss(\params)=\tfrac12\ipS{\W}{\W}-\ipS{\W}{\Wt}+\tfrac12\ipS{\Wt}{\Wt},
\end{equation}
and every term is an exact contraction, with no sampling noise:
\begin{itemize}
  \item For two TTNs $\W,\W'$ on the same tree (bond dimensions may differ), environments
    $E_v\in\R^{\bd{v}\times\bd{v}'}$ are computed from the leaves to the root,
    \begin{equation}\label{eq:environment}
      (E_v)_{aa'}=\sum\G{v}\big[(j_i),(b_c),a\big]\,\G{v}'\big[(j'_i),(b'_c),a'\big]
      \prod_{i\in\sites{v}}(\Sigi{i})_{j_ij'_i}\prod_{c\in\ch{v}}(E_c)_{b_cb'_c},
    \end{equation}
    and $\ipS{\W}{\W'}=E_{\vroot}$. The cost is linear in $\Kc$ at fixed $d$ and $\chi$.
  \item For a CP teacher, $\ipS{\W}{\U{r}}$ is the student evaluated at the site vectors
    $\Sigi{i}\fac{r}{i}$, so for any tree
    \begin{equation}\label{eq:cp-inner}
      \ipS{\W}{\Wt}=\sum_r\tw{r}\,f_\params\big(\Sigi{1}\fac{r}{1},\dots,\Sigi{\Ns}\fac{r}{\Ns}\big),
      \qquad
      \ipS{\Wt}{\Wt}=\sum_{r,r'}\tw{r}\tw{r'}\prod_i\fac{r}{i}{}\T\Sigi{i}\fac{r'}{i},
    \end{equation}
    where $f_\params(z_1,\dots,z_\Ns)$ denotes \eqref{eq:messages} with $\feat{i}(x)$ replaced by
    $z_i$.
  \item A TTN teacher on another tree is glued to the student along the physical legs; the
    resulting network has loops and is contracted exactly along an optimized contraction path.
\end{itemize}
For an odeco teacher with whitened inputs, \eqref{eq:loss-expanded} reads
$\Loss=\tfrac12\normF{\W}^2-\sum_r\tw{r}^2\ov{r}+\tfrac12\sum_r\tw{r}^2$. A \emph{partial fit}
$\W_S=\sum_{r\in S}\tw{r}\U{r}$ that has learned exactly the components in $S$ has loss
$\tfrac12\sum_{r\notin S}\tw{r}^2$; in particular, learning the components in decreasing order of
$\tw{r}$ passes through the levels
\begin{equation}\label{eq:partial-fit}
  \Lpart{k}=\tfrac12\sum_{r>k}\tw{r}^2,\qquad k=0,\dots,\ncomp.
\end{equation}
These are the candidate plateau levels examined in \cref{sec:s2s}.

\paragraph{Anisotropic and correlated inputs.}
Anisotropic product covariances are handled exactly by \eqref{eq:loss}--\eqref{eq:cp-inner}; the
default family has power-law site spectra,
$\Sigi{i}=c_\beta\diag(1^{-\beta},2^{-\beta},\dots,d^{-\beta})$ with $c_\beta$ fixed by
$\tr\Sigi{i}=d$, aligned with the teacher's factors or rotated away from them. Cross-site
correlations break the product structure of $\Sig$; they are studied with sampled data only
(\cref{sec:covariance-teacher}).

\paragraph{Finite samples.}
For sample-size questions, a training set $(x_j,y_j)_{j=1}^{\nsamp}$ has
$x_j\sim\Normal(0,\Sig)$ \iid and $y_j=f^*(x_j)+\varepsilon_j$ with optional label noise
$\varepsilon_j\sim\Normal(0,\varsigma^2)$. The empirical loss is
$\Lemp(\params)=\frac{1}{2\nsamp}\sum_{j}(f_\params(x_j)-y_j)^2$, and the test error is the
population loss \eqref{eq:loss}.

% ===========================================================================================
\subsection{Initialization and its variants}\label{sec:setup-init}

\paragraph{Fan-in initialization.}
The default initialization is
\begin{equation}\label{eq:init}
  \G{v}=\frac{\iscv{v}}{\sqrt{\fanin{v}}}\,Z_v,\qquad Z_v \text{ with \iid } \Normal(0,1)
  \text{ entries},
\end{equation}
with a single per-core scale $\iscv{v}=\isc$ unless stated otherwise. Each core is then
approximately $\iscv{v}$ times an isometry (when $\fanin{v}\gg\bd{v}$), so that with $\iscv{v}=1$
the messages \eqref{eq:messages} keep their size on the way up the tree.

\begin{proposition}[Initial scales]\label{prop:init-norm}
  \proven{} Under \eqref{eq:init}, $\E\normF{\W(\params_0)}^2=\prod_{v}\iscv{v}^2$ exactly,
  $\E\normF{\G{v}}^2=\iscv{v}^2\bd{v}$, and hence $\E\norm{\params_0}^2=\sum_v\iscv{v}^2\bd{v}$.
\end{proposition}
\begin{proof}
  Let $\Sub{v}$ be the contraction of the subtree below $v$ with its parent leg open, a matrix with
  $\bd{v}$ columns. Conditionally on the children's subtrees, the independence and variance
  $\iscv{v}^2/\fanin{v}$ of the entries of $\G{v}$ give
  $\E[\Sub{v}\T\Sub{v}]=\frac{\iscv{v}^2}{\fanin{v}}\prod_{i\in\sites{v}}\pd[i]
  \prod_{c\in\ch{v}}\tr(\Sub{c}\T\Sub{c})\,\Id_{\bd{v}}$. The children's subtrees are independent, so
  induction from the leaves gives $\E[\Sub{v}\T\Sub{v}]=\big(\prod_{u}\iscv{u}^2\big)\Id_{\bd{v}}$,
  the product running over the subtree of $v$, because $\fanin{v}=\prod_{i\in\sites{v}}\pd[i]
  \prod_{c}\bd{c}$. At the root $\bd{\vroot}=1$ and $\Sub{\vroot}$ is $\W$ flattened. The second
  claim is $\E\normF{\G{v}}^2=\fanin{v}\bd{v}\,\iscv{v}^2/\fanin{v}$.
\end{proof}

We therefore report both the per-core scale $\isc$ and the function-space scale
$\fsc=\isc^{\Kc}$, the root-mean-square norm of $\W(\params_0)$. Small $\isc$ makes $\W(\params_0)$
exponentially small in $\Kc$, while $\norm{\params_0}^2\approx\isc^2\sum_v\bd{v}$ grows with the
bond dimensions.

\paragraph{Variants.}
\begin{enumerate}[label=(\alph*)]
  \item \emph{Imbalance}: per-core scales $\iscv{v}=\isc c_v$ with a fixed profile $(c_v)$, for
    example root versus the rest, or by depth, compared at fixed function-space scale
    $\prod_v\iscv{v}$.
  \item \emph{Centred model}: the target is $\Wt+\W(\params_0)$, that is, $f_\params-f_{\params_0}$ is
    trained on $f^*$ \citep{chizat2019lazy}. The initial function is removed exactly at any $\isc$
    without changing the parameterization, so $\isc$ alone moves the model from small to large
    initialization. Overlaps are measured on $\W-\W(\params_0)$. This is the main construction of the
    regime map (\cref{sec:regime-map}).
  \item \emph{Two-copy model} \citep{woodworth2020kernel}: $\W(\params^+)-\W(\params^-)$ with
    $\params^+(0)=\params^-(0)=\params_0$, which vanishes at initialization for every $\isc$. It is
    represented as one TTN on the same tree with doubled bond dimensions $2\bd{v}$ and block-diagonal
    cores: core $v$ holds $\G{v}^+$ on the block where all its bond indices are in the ``$+$'' half and
    $\G{v}^-$ on the ``$-$'' block, except that the root holds $\G{\vroot}^+$ and
    $-\G{\vroot}^-$. It is used as a cross-check of the centred model; see \cref{rem:two-copy-mask}.
  \item \emph{Aligned (decoupled) initialization} for CP teachers (\cref{prop:aligned}): cores
    diagonal in the teacher's basis, giving exactly solvable Saxe-type dynamics.
  \item \emph{Per-core learning-rate multipliers} $\lrm{v}$ (\cref{sec:setup-dynamics}). A fixed
    forward multiplier $c_v$ on core $v$ (writing $\G{v}=c_vH_v$ and training $H_v$) is equivalent to
    the initial scale $c_v\iscv{v}$ together with the multiplier $c_v^2\lrm{v}$. This is how NTK-type
    and $\mu$P-type parameterizations for large bond dimension \citep{yang2021tensor} are expressed
    (\cref{sec:other-regimes}).
\end{enumerate}

\begin{remark}[The two-copy model needs a gradient mask]\label{rem:two-copy-mask}
  \proven{} Gradient flow on the doubled TTN does not preserve its block-diagonal structure. For a
  root with children $c_1,c_2$, the gradient of $\Loss$ with respect to the root entry pairing index
  $b_1$ of the $+$ half of bond $\edge{c_1}$ with index $b_2$ of the $-$ half of bond $\edge{c_2}$ is
  $\ipS{\W-\Wt}{\Sub{c_1}^+[:,b_1]\otimes\Sub{c_2}^-[:,b_2]}$ (up to the ordering of sites), which at
  initialization equals $-\ipS{\Wt}{\Sub{c_1}^+[:,b_1]\otimes\Sub{c_2}^-[:,b_2]}$ and is generically
  nonzero; mixed blocks of non-root cores also receive gradient. The two-copy dynamics is therefore
  the masked flow $\dot\params=-\mask\odot\nabla\Loss(\params)$, with $\mask$ the $0/1$ indicator of
  the diagonal blocks, which keeps the blocks and equals gradient flow of
  $\tfrac12\normS{\W(\params^+)-\W(\params^-)-\Wt}^2$ in $(\params^+,\params^-)$.
\end{remark}

\begin{proposition}[Aligned initialization]\label{prop:aligned}
  \proven{} \knownresult[deep linear networks: \citealp{saxe2014exact}] Let $\Wt$ be odeco, let
  inputs be whitened, let $\bd{v}\ge\ncomp$ for every edge, and let every subtree and every
  subtree's complement carry at least one site (true for the trees of \cref{tab:trees}).
  Initialize
  \begin{equation}\label{eq:aligned-init}
    \G{v}=\sum_{r=1}^{\ncomp}\amp{v}{r}\Big(\bigotimes_{i\in\sites{v}}\fac{r}{i}\Big)\otimes
    \Big(\bigotimes_{c\in\ch{v}}\varepsilon_r\Big)\otimes\varepsilon_{r(v)},
  \end{equation}
  with $\varepsilon_r$ the $r$-th standard basis vector of the bond in question, $r(v)=r$ for
  $v\ne\vroot$ and $r(\vroot)=1$. Under gradient flow with multipliers $\lrm{v}$, the cores keep the
  form \eqref{eq:aligned-init}, $\W=\sum_r\prodamp{r}\U{r}$ with
  $\prodamp{r}=\prod_v\amp{v}{r}$, $\ov{r}=\prodamp{r}/\tw{r}$,
  $\Loss=\tfrac12\sum_r(\tw{r}-\prodamp{r})^2$, and
  \begin{equation}\label{eq:aligned-ode}
    \dot q_{v,r}=\lrm{v}\,(\tw{r}-\prodamp{r})\prod_{u\ne v}\amp{u}{r}.
  \end{equation}
  With $\lrm{v}=1$ the differences $\amp{v}{r}^2-\amp{u}{r}^2$ are conserved; a balanced start
  $\amp{v}{r}(0)=q_r(0)$ stays balanced, and $\dot q_r=q_r^{\Kc-1}(\tw{r}-q_r^{\Kc})$.
\end{proposition}
\begin{proof}
  On the family \eqref{eq:aligned-init}, contracting the basis vectors on the bonds gives
  $\W=\sum_r\prodamp{r}\U{r}$, and the $\U{r}$ are orthonormal. The derivative of $\W$ with respect
  to an entry of $\G{v}$ is the tensor product of basis vectors on the physical legs of $v$, the
  columns $\Sub{c}[:,b_c]$ of the child subtrees and the column of the environment above $v$
  indexed by the parent index $a$. On the family these columns are multiples of
  $\bigotimes_{i\in\below{c}}\fac{b_c}{i}$ (and of the analogous product over the sites outside
  $\below{v}$), so the inner product with $\W-\Wt=\sum_r(\prodamp{r}-\tw{r})\U{r}$ vanishes unless all
  bond indices equal the same $r$ (with $a=1$ at the root). The gradient therefore lies in the span
  of the family, and its component along $\amp{v}{r}$ is
  $(\prodamp{r}-\tw{r})\prod_{u\ne v}\amp{u}{r}$. Finally, with $\lrm{v}=1$,
  $\ddt\amp{v}{r}^2=2(\tw{r}-\prodamp{r})\prodamp{r}$ is the same for every $v$.
\end{proof}

For $\Kc\ge3$ a balanced mode started at $q_r(0)\ll\tw{r}^{1/\Kc}$ reaches $\ov{r}=\frac12$ at
\begin{equation}\label{eq:aligned-escape}
  t_r=\frac{q_r(0)^{-(\Kc-2)}}{(\Kc-2)\,\tw{r}}+O\big(\tw{r}^{-(2-2/\Kc)}\big),
\end{equation}
and for $\Kc=2$ at $t_r=\log(\tw{r}/q_r(0)^2)/(2\tw{r})+O(q_r(0)^2/\tw{r}^2)$
(direct integration of the balanced ODE) \proven. With $q_r(0)\propto\isc$ the plateau before
the first escape grows like $\isc^{-(\Kc-2)}$; this time scale is why the integration below must be
adaptive and in double precision.

% ===========================================================================================
\subsection{Training dynamics}\label{sec:setup-dynamics}

\paragraph{Gradient flow} is the reference dynamics,
\begin{equation}\label{eq:gf}
  \dot\params=-\Pre\,\nabla\Loss(\params),
\end{equation}
where $\Pre$ is diagonal with entry $\lrm{v}\mask_j$ for every parameter $j$ of core $v$: the
$\lrm{v}$ are the per-core learning-rate multipliers (default $1$) and $\mask\in\{0,1\}^{\dim\params}$
is an optional mask (\cref{rem:two-copy-mask}; default all ones). Here $t$ always denotes
gradient-flow time. It is integrated with the adaptive Dormand--Prince 5(4) pair
\citep{dormand1980family} (first-same-as-last, local extrapolation) in double precision, with
relative and absolute tolerances $10^{-8}$ and $10^{-12}$ by default. A step of size $h$ is accepted
when the root-mean-square over components of the embedded error estimate, each divided by
$10^{-12}+10^{-8}\max(\abs{\params_j},\abs{\params_j^{\mathrm{new}}})$, is at most one; the next step is
$h\min\{5,\max\{0.2,0.9\,\mathrm{err}^{-1/5}\}\}$, with no increase right after a rejection, and the
first step follows \citet[Sec.~II.4]{hairer1993solving}. Step sizes thus follow the local time scale
$\norm{\params}/\norm{\dot\params}$, so that long plateaus are cheap. Diagnostics are recorded at
$t=0$ and at $200$ log-spaced times from $t_{\min}=10^{-3}\min(\tau_0,\tmax)$ to $\tmax$, where
$\tau_0=\norm{\params_0}/\norm{\dot\params(0)}$; steps are shortened to land on these times exactly.
A run stops at $\tmax$, when $\Loss$ falls below a prescribed floor, after a maximum number of steps,
or if the step size underflows. Accuracy is checked with the conservation laws
(\cref{sec:setup-conservation}) and against the closed-form aligned solutions (\cref{prop:aligned}).

\paragraph{Gradient descent} with learning rate $\lr$,
$\params_{k+1}=\params_k-\lr\Pre\nabla\Loss(\params_k)$, is recorded at $t=\lr k$, so it approaches
gradient flow at equal $t$ as $\lr\to0$. A run has diverged when $\Loss$ becomes non-finite or
exceeds $10^{6}\Loss(\params_0)$.

\paragraph{Stochastic gradient descent} uses the minibatch loss
$\hat\Loss_{\batch}(\params)=\frac{1}{2\batch}\sum_{b=1}^{\batch}(f_\params(x_b)-y_b)^2$ in the
same update, again recorded at $t=\lr k$. Online SGD draws fresh inputs $x_b\sim\Normal(0,\Sig)$
with $y_b=f^*(x_b)$ at every step; with a fixed training set, batches are taken from a new random
permutation in every epoch. Optional label noise is drawn afresh at every step.

\paragraph{Precision.}
All computations are in double precision: with $\isc=10^{-2}$ and $\Kc=7$,
$\normF{\W(\params_0)}\approx10^{-14}$ and the gradients are of order $10^{-12}$, which single
precision cannot resolve over long plateaus.

% ===========================================================================================
\subsection{Symmetries and conservation laws}\label{sec:setup-conservation}

For a bond $e=\edge{v}$ joining core $v$ to its parent $p=\pa{v}$, let $\Gm{u}{e}$ denote core
$u\in\{v,p\}$ matricized with its $e$-leg as rows (a $\bd{v}\times\fanin{v}$ matrix for $u=v$). For
every invertible $B\in\GL(\bd{v})$ the \emph{gauge transformation}
\begin{equation}\label{eq:gauge}
  \Gm{v}{e}\mapsto B\,\Gm{v}{e},\qquad \Gm{p}{e}\mapsto B^{-\top}\Gm{p}{e}
\end{equation}
leaves $\W$, and hence any loss that depends on $\params$ only through $\W$, unchanged.

\begin{proposition}[Conservation laws]\label{prop:conservation}
  \proven{} \knownresult[deep linear networks: \citealp{arora2018optimization,du2018algorithmic}]
  Under gradient flow $\dot{\G{u}}=-\lrm{u}\nabla_{\G{u}}\Loss$ of any loss that depends on $\params$
  only through $\W$ (population or empirical), the $\bd{v}\times\bd{v}$ matrix
  \begin{equation}\label{eq:conserved}
    \Q{e}=\frac{1}{\lrm{v}}\Gm{v}{e}\Gm{v}{e}\T-\frac{1}{\lrm{p}}\Gm{p}{e}\Gm{p}{e}\T
  \end{equation}
  is constant in time, for every bond $e=\edge{v}$ with $p=\pa{v}$.
\end{proposition}
\begin{proof}
  Write $X=\Gm{v}{e}$, $Y=\Gm{p}{e}$ and take $B=\Id+\epsilon M$ in \eqref{eq:gauge}, so that
  $X\mapsto X+\epsilon MX$ and $Y\mapsto Y-\epsilon M\T Y+O(\epsilon^2)$. Invariance of $\Loss$ at
  first order in $\epsilon$ gives $\ip{\nabla_X\Loss\,X\T-Y\,(\nabla_Y\Loss)\T}{M}=0$ for every $M$,
  so $C:=\nabla_X\Loss\,X\T=Y(\nabla_Y\Loss)\T$. Along the flow,
  $\ddt(XX\T)/\lrm{v}=-(C+C\T)=\ddt(YY\T)/\lrm{p}$.
\end{proof}

This is the tree version of the balancedness of deep linear networks and an instance of the
general correspondence between continuous symmetries and conserved quantities of gradient flow
\citep{kunin2021neural}. Consequences:
\begin{itemize}
  \item \emph{Norm balance}: $\tr\Q{e}=\normF{\G{v}}^2/\lrm{v}-\normF{\G{p}}^2/\lrm{p}$. With equal
    multipliers, the differences $\normF{\G{u}}^2-\normF{\G{u'}}^2$ between any two cores are
    conserved (the tree is connected), so all cores grow together.
  \item \emph{Balancedness at small initialization}: under \eqref{eq:init},
    $\E\,\Q{e}(0)=\big(\iscv{v}^2/\lrm{v}-\iscv{p}^2\bd{p}/(\lrm{p}\bd{v})\big)\Id$, so with small
    $\isc$ the cores remain balanced up to $O(\isc^2)$ while they grow to $O(1)$. Imbalanced
    initializations or multipliers set $\Q{e}$ to other values, a second knob for the regime map
    \citep{kunin2024get}.
  \item The masked two-copy flow conserves the $\Q{e}$ of each copy (its gauge group is the
    block-diagonal subgroup). Gradient descent and SGD conserve $\Q{e}$ only up to $O(\lr^2)$ per step.
  \item \emph{Integrator check}: the conservation drift
    \begin{equation}\label{eq:drift}
      \mathrm{drift}(t)=\frac{\max_e\normmax{\Q{e}(t)-\Q{e}(0)}}
      {\max_e\normmax{\Q{e}(0)}+\max_u\normF{\G{u}(t)}^2/\lrm{u}},
    \end{equation}
    with $\normmax{\cdot}$ the largest absolute entry, measures the integration error relative to the
    size of the terms that cancel in $\Q{e}$. It must stay below $10^{-6}$ under gradient flow.
\end{itemize}

% ===========================================================================================
\subsection{Regime diagnostics}\label{sec:setup-diagnostics}

Regimes are assigned from measured quantities, never by eye. The thresholds below are conventions;
every phase diagram reports how its boundaries move when they change.

\paragraph{Raw quantities} (\cref{tab:raw-diagnostics}) are recorded on the logarithmic time grid of
\cref{sec:setup-dynamics}, with extra points where needed. The tangent kernel uses a fixed probe set
of $\nprobe=128$ inputs $x_1,\dots,x_{\nprobe}\sim\Normal(0,\Sig)$ and the Jacobian
$\Jac_t=\partial\big(f_\params(x_1),\dots,f_\params(x_{\nprobe})\big)/\partial\params$ at
$\params(t)$. The gauge directions at $\params$ are the tangent vectors of the orbits
\eqref{eq:gauge}, $(\delta\Gm{v}{e},\delta\Gm{p}{e})=(M\Gm{v}{e},-M\T\Gm{p}{e})$ for all bonds $e$
and $M\in\R^{\bd{v}\times\bd{v}}$; $\Hperp$ is the Hessian restricted to their orthogonal
complement.
% TODO(regime-map): with multipliers or a mask, decide whether the kernel is J J^T or J D J^T
% (the latter drives the function-space dynamics of eq:gf) and state it here.

\begin{table}[t]
  \centering
  \caption{Raw diagnostics recorded along every run.}
  \label{tab:raw-diagnostics}
  \begin{tabularx}{\textwidth}{@{}lX@{}}
    \toprule
    symbol & definition \\
    \midrule
    $\Loss(t)$, $\lossn(t)$ & exact population loss \eqref{eq:loss}; normalized loss
      $\lossn(t)=\Loss(t)/\Loss(0)$ \\
    $\bs{e}{k}(t)$ & bond spectrum of $\W(t)$ across edge $e$; function-space version when
      $\Sig\ne\Id$ \\
    $\bst{e}{k}$ & bond spectrum of $\Wt$; equal to the weights $\tw{k}$ on every edge for an odeco
      teacher (\cref{prop:teacher-spectra}) \\
    $\ov{r}(t)$ & overlap \eqref{eq:overlap} with teacher component $r$ (CP teachers); of
      $\W-\W(\params_0)$ in the centred model \\
    $\dth(t)$ & $\norm{\params(t)-\params(0)}/\norm{\params(0)}$, also per core \\
    $\TK_t$ & tangent kernel $\Jac_t\Jac_t\T$ on the probe set; for small $\Ns$ also the
      tensor-space operator $\sum_v(\partial\W/\partial\G{v})(\partial\W/\partial\G{v})\T$ \\
    $\dTK(t)$ & relative kernel change $\normF{\TK_t-\TK_0}/\normF{\TK_0}$ \\
    $\KTA(t)$ & kernel-target alignment $y\T\TK_ty/(\normF{\TK_t}\norm{y}^2)$, with
      $y_p=f^*(x_p)$ on the probe set \\
    $\mathrm{drift}(t)$ & conservation drift \eqref{eq:drift}, an integrator check \\
    $\Hess(t)$, $\Hperp(t)$ & Hessian of $\Loss$ at $\params(t)$, and its restriction to the
      complement of the gauge directions \\
    \bottomrule
  \end{tabularx}
\end{table}

\paragraph{Learning events.}
For a CP teacher, event $r$ occurs at the first time $\tev{r}$ at which $\ov{r}(t)$ crosses $\frac12$,
and its width is $\wev{r}=t(\ov{r}=0.9)-t(\ov{r}=0.1)$. For a general teacher, event $(e,k)$ occurs
when $\bs{e}{k}(t)$ first crosses half of its final value, with widths defined in the same way from
the crossings of $10\%$ and $90\%$ of the final value; since one component typically raises several
bond spectra at once, events closer than their widths are merged.
% TODO(diagnostics): notes/diagnostics.md does not define the width of spectral events; the
% 10%-90% rule above is the natural analogue of the CP rule. Confirm.
The \emph{effective bond rank} is $\rk{e}(t)=\#\{k:\bs{e}{k}(t)>0.1\,\bst{e}{\min}\}$, where
$\bst{e}{\min}$ is the smallest nonzero teacher Schmidt value on $e$; the threshold-free alternative is
the entropy rank $\rkH{e}=\exp(-\sum_kp_k\log p_k)$ with
$p_k=(\bs{e}{k})^2/\sum_j(\bs{e}{j})^2$.

\paragraph{Plateaus and the saddle-to-saddle signature.}
Number the events in time order, $\tev{1}<\tev{2}<\cdots$, with widths $\wev{i}$, and define the
\emph{separation ratios}
\begin{equation}\label{eq:separation}
  \sep{i}=\frac{\tev{i+1}-\tev{i}}{\max(\wev{i},\wev{i+1})}.
\end{equation}
Between two events with $\sep{i}>3$ the loss sits on a plateau at the level of a partial fit. The
plateau midpoint $\tmid{i}$ is the time of minimal speed $\norm{\dot\params(t)}$ in
$[\tev{i},\tev{i+1}]$, and the Hessian is evaluated there. A run shows the
\emph{saddle-to-saddle (S2S) signature} when
\begin{enumerate}[label=(\roman*)]
  \item there are at least two events,
  \item every separation ratio satisfies $\sep{i}>3$, and
  \item across a sweep, plateau durations grow without bound as the initialization scale tends to
    zero (like $\log(1/\isc)$ or a power of $1/\isc$).
\end{enumerate}
Plateaus are judged in linear time between events, not only on logarithmic time axes: an escape that
needs two cores takes a time of order $\log(1/\isc)$, which can be invisible on a log-time axis after
a long first plateau.

\paragraph{Regime assignment.}
Each run, trained up to time $T$, receives one label by the rules of \cref{tab:regimes}.

\begin{table}[t]
  \centering
  \caption{Regime assignment for one run trained up to time $T$ (the S2S signature is defined in
    \cref{sec:setup-diagnostics}). Threshold sensitivity: the kernel threshold $0.1$ is varied over
    $\{0.05,0.1,0.2\}$ and the separation threshold $3$ over $\{2,3,5\}$.}
  \label{tab:regimes}
  \begin{tabularx}{\textwidth}{@{}lX@{}}
    \toprule
    regime & rule \\
    \midrule
    lazy & the loss converges and $\dTK(T)<0.1$ (so $\dth$ is small as well) \\
    saddle-to-saddle & $\dTK(T)\ge0.1$ and the S2S signature (i)--(iii) \\
    rich, non-incremental & $\dTK(T)\ge0.1$ and the loss converges, but some $\sep{i}\le3$
      (overlapping transitions) \\
    stalled & $\Loss(T)\ge10^{-3}\Loss(0)$ within the time budget (for example a long initial
      plateau) \\
    diverged & GD and SGD only: $\Loss$ non-finite or above $10^{6}\Loss(0)$ \\
    \bottomrule
  \end{tabularx}
\end{table}

\begin{remark}[Reading of the rules]\label{rem:regime-rules}
  ``The loss converges'' means $\Loss(T)<10^{-3}\Loss(0)$, the complement of ``stalled''. Labels are
  checked in the order diverged, stalled, then the three converged regimes. Condition (iii) is a
  property of a sweep: a single run is labelled S2S on the basis of (i)--(ii), and a region of a
  phase diagram is reported as S2S only if (iii) holds there. A converged run with $\dTK(T)\ge0.1$
  and a single (merged) event is counted as rich, non-incremental.
  % TODO(lead): these four readings make the rules of notes/diagnostics.md exhaustive; confirm or
  % amend them in notes/diagnostics.md.
\end{remark}

\paragraph{Saddle type at a plateau.}
At a plateau midpoint, a \emph{strict saddle} has $\lambda_{\min}(\Hperp)<-\epsilon$ with
$\epsilon=10^{-3}\max_j\abs{\lambda_j(\Hperp)}$; a \emph{higher-order (non-strict) saddle} has no
eigenvalue below $-\epsilon$ although the loss later decreases. We report the number of negative
eigenvalues, $\lambda_{\min}$ (also unnormalized) and $\lambda_{\min}/\lambda_{\max}$.

\paragraph{Reporting rules.}
Every scaling claim uses at least five initialization scales spanning at least one decade and at
least five seeds per point, reports mean $\pm$ standard error and a fitted exponent with its standard
error, and overlays the theoretical prediction. Every regime boundary comes with the threshold
sensitivity check of \cref{tab:regimes}. Every number states the dynamics that produced it
(gradient flow with its tolerances, GD with its step size, or SGD with its step size and batch size).

% ===========================================================================================
\subsection{What the results may depend on}\label{sec:setup-scope}

\begin{itemize}
  \item \emph{Population versus finite samples.} Plateaus and saddle sequences are properties of the
    population gradient flow; statements about implicit bias and generalization need finite samples
    and are tested separately.
  \item \emph{Multilinear features.} There is no pointwise nonlinearity; orthonormal local feature
    maps all give the same population dynamics, and non-orthonormal ones are equivalent to anisotropic
    $\Sigi{i}$.
  \item \emph{Product input measure.} It makes the population loss a tree contraction; cross-site
    correlations break this and are studied with sampled data.
  \item \emph{Placement of physical legs.} Binary trees (two sites per bottom core) and HT (one leaf
    core per site) differ in $\Kc$ and in how many cores must change to learn a component, which is
    expected to change plateau-time exponents.
  \item \emph{Scalar output.} Multi-output regression or classification would add an output leg to
    the root core.
  \item \emph{Fan-in initialization.} It fixes the relative scales of the cores. Exponents should
    not depend on it, but prefactors and the order of learning might, which is why imbalanced and
    aligned initializations are also studied.
\end{itemize}
